% Sections 3.1-3.9, from lectures/L03/appendix/a_flip_flops_and_registers.md (A.1-A.9).

\section{From combinational to sequential networks}
\label{c3:sec:seq}

Every circuit in \chapref{c1:ch} and \ref{c2:ch} was \textbf{combinational}: its output is a pure
function of its current inputs, so the same inputs always give the same output, and there is no
notion of ``before'' or ``after'' inside the circuit.

Most useful systems need more. A counter has to remember its count. A register has to hold a value
after the signal that produced it has changed. An LED toggled by a button has to remember whether
it is lit, since that is not encoded in the button's momentary position at all.

All of these need \textbf{memory}: an output that depends on the circuit's own \emph{previous}
output. A sequential network is combinational logic with part of the output fed back to the input,
and that \textbf{feedback loop} is what turns a stateless gate network into a circuit with state.
This chapter builds the primitive that makes controlled feedback practical, the D flip-flop, and
then the two things you build directly from it: registers and edge detectors.

% ------------------------------------------------------------------------------------------
\section{The D latch}
\label{c3:sec:latch}

The simplest storage element is the \textbf{D latch}: two cross-coupled gates, gated by an
\code{enable} input, storing a single bit.

\begin{textdrawing}
        ┌────────────┐
   D ──►│            │──► Q
        │  D LATCH   │
enable─►│            │──► Qn
        └────────────┘
\end{textdrawing}

\code{D} is the value to store, \code{enable} controls whether the latch is \emph{open}
(transparent) or \emph{closed} (holding), \code{Q} is the stored value and \code{Qn} is always its
inverse.

Each output feeds back into the other's equation, which is exactly \secref{c3:sec:seq}'s feedback
loop:

\[ Q = (D' \cdot \mathit{enable} + Q_n)' \]
\[ Q_n = (D \cdot \mathit{enable} + Q)' \]

\begin{booktable}{@{}p{4em}p{12em}L@{}}
  \thead{\code{enable}} & \thead{Behaviour} & \thead{Result} \\
  \midrule
  \code{1} & the latch is \textbf{open} (transparent) & $Q = D$, $Q_n = D'$; the output follows the
    input immediately \\
  \code{0} & the latch is \textbf{closed} (latched) & \code{Q} and \code{Qn} hold whatever they
    last had \\
\end{booktable}

\textbf{Build it in CircuitVerse} from a NOT gate, two AND gates and two NOR gates (a NOR is an OR
followed by a NOT). Cross-couple the feedback so that the NOR producing \code{Q} feeds the NOR on
the \code{Qn} side and the other way round. Drive \code{D} and \code{enable} from input switches
and watch \code{Q} and \code{Qn} on LEDs.

Then test it:
\begin{itemize}
  \item With \code{enable = 1}, change \code{D} and confirm that \code{Q} follows immediately. The
    latch is \textbf{transparent}.
  \item With \code{enable = 0}, confirm that \code{Q} is unchanged whatever \code{D} does. The
    latch is \textbf{latched}.
  \item Toggle \code{D} a few times while switching \code{enable}, and confirm that the latch
    always freezes whatever \code{Q} held at the moment \code{enable} went low. That is what
    \textbf{level-sensitive} means: transparent for the whole time \code{enable} is high, not just
    at a single instant.
\end{itemize}

That last point is also the latch's fundamental weakness. Because it is transparent for the whole
time \code{enable = 1}, every glitch on \code{D} during that window passes straight through to
\code{Q}. In a large synchronous system, guaranteeing stable data for a whole clock phase is far
harder than guaranteeing it around a single instant. That is the problem the D flip-flop solves.

% ------------------------------------------------------------------------------------------
\section{The D flip-flop}
\label{c3:sec:dff}

A \textbf{D flip-flop} stores one bit the way a latch does, but is \textbf{edge-triggered} rather
than level-sensitive: it samples \code{D} only at the \emph{instant} the clock changes, not for the
whole time the clock is high or low.

Feed both the same \code{D} and the same clock, and the difference shows immediately:

\lecfig[0.78]{lectures/L03/appendix/images/latch_vs_flipflop.png}{A latch follows \code{D} during
every window where the clock is high, while a flip-flop changes only at the rising
edges.}{c3:fig:latchvsff}

The latch changes three times: twice because \code{D} moved while the clock was high, and once
because the clock went high against a \code{D} that had already changed. The flip-flop changes
once, because only one of the four rising edges found \code{D} different from the value \code{Q}
already held. Both behave correctly; they answer different questions.

\begin{textdrawing}
        ┌──────────────────┐
   D ──►│                  │──► Q
clock ─►│   D FLIP-FLOP    │
enable─►│  (rising edge)   │──► Qn
reset_n►│                  │
        └──────────────────┘
\end{textdrawing}

\code{clock} is the time reference (\secref{c3:sec:clock}). \code{enable} gates whether the
flip-flop takes a new value \emph{at the triggering edge}. \code{reset_n} is an asynchronous,
active-low reset: whenever it is \code{0}, \code{Q} is forced to \code{0} immediately, regardless
of the clock.

\begin{booktable}{@{}p{15em}L@{}}
  \thead{Condition} & \thead{Result} \\
  \midrule
  \code{reset_n = 0} & $Q = 0$, $Q_n = 1$, immediately (asynchronous, ignoring the clock
    entirely) \\
  \code{reset_n = 1}, rising edge on \code{clock}, \code{enable = 1} & $Q = D$, $Q_n = D'$, sampled
    at that instant \\
  \code{reset_n = 1}, rising edge on \code{clock}, \code{enable = 0} & \code{Q} holds its previous
    value \\
  \code{reset_n = 1}, no rising edge on \code{clock} & \code{Q} holds its previous value, whatever
    \code{D} does \\
\end{booktable}

That last row is worth seeing rather than reading:

\lecfig[0.72]{lectures/L03/appendix/images/dff_timing.png}{\code{D} goes high mid-cycle and is
caught at the next rising edge, then dips and recovers entirely between two edges without \code{Q}
ever moving.}{c3:fig:dfftiming}

\code{D} dips low and comes back within a single clock period, and \code{Q} never moves, because no
edge occurred while it was low. A latch in the same position would have followed it down and up
again.

Internally, an edge-triggered flip-flop is two D latches in series, a \textbf{master-slave} pair.
One is transparent while the clock is low and captures \code{D}; the other is transparent while the
clock is high and passes the captured value on to \code{Q}. The net effect is that \code{Q} changes
only at the moment the clock goes from low to high, never while it sits at either level. The
external behaviour is what VHDL's \code{rising_edge()} models (\secref{c3:sec:template}).

\textbf{Build it in CircuitVerse} by building on your latch from \secref{c3:sec:latch}:
\begin{itemize}
  \item Duplicate the latch so that two sit in series: the \emph{master} feeds the \emph{slave},
    whose \code{Q} is the flip-flop's output.
  \item Add \code{clock}, \code{reset_n}, \code{D} and \code{enable} as inputs.
  \item Drive the two latches' transparency from the clock and its inverse. To keep the two
    meanings apart, call the internal per-latch control signal \code{gate}: the master's
    \code{gate} is \code{clock'} and the slave's is \code{clock}. The flip-flop's external
    \code{enable} is a different signal and is connected separately, below.
  \item Connect the external \code{enable} as a 2-to-1 mux (\secref{c2:sec:mux}) on the master's
    data input: when \code{enable = 1} the mux passes \code{D} through, and when \code{enable = 0}
    it passes the flip-flop's own \code{Q} back in, so that the edge recaptures the value already
    stored. Gate the \emph{data}, never the clock.
  \item Add the reset to \textbf{both} latches. Forcing \code{Q} low takes two changes in each
    latch, not one:
  \begin{itemize}
    \item feed an active-high \code{reset} (\code{reset_n} through a NOT) into the NOR producing
      \code{Q}, which forces \code{Q} to \code{0}.
    \item AND \code{reset_n} into the term \code{D AND gate} in the same latch, which frees
      \code{Qn} to go high.
  \end{itemize}
  \item Set the clock period to something slow enough to watch, e.g. \code{1000 ms}.
\end{itemize}

With the reset wired in, the equations for each latch become:

\[ Q = (D' \cdot \mathit{gate} + Q_n + \mathit{reset})' \]
\[ Q_n = (D \cdot \mathit{gate} \cdot \mathit{reset\_n} + Q)' \]

A tempting shortcut is to leave the \code{Qn} equation alone and break the loop instead, by ANDing
\code{reset_n} into the \code{Qn} wire that feeds \code{Q}'s NOR. That is not enough. With
\code{gate = 1} and \code{D = 1} it leaves \code{Qn} at \code{0}, which is what holds \code{Q} at
\code{1}, so the reset silently does nothing in exactly the case you are most likely to test it in.
Both changes, in both latches, or neither.

Confirm that \code{Q} updates only at the moment the clock goes high, and sits still between edges
even if \code{D} changes, and that a low \code{reset_n} forces \code{Q} to \code{0} immediately,
while releasing it changes nothing until the next rising edge.

% ------------------------------------------------------------------------------------------
\section{Clocking}
\label{c3:sec:clock}

The \textbf{clock} is a periodic square wave, which every flip-flop in a synchronous circuit reads
from the same source. Two numbers describe it: the \textbf{period} $T$, the time for one whole
cycle, and the \textbf{frequency} $f = 1/T$ in Hz.

Every period has exactly two \textbf{edges}: a \textbf{rising edge} from \code{0} to \code{1}, and
a \textbf{falling edge} from \code{1} to \code{0}. A synchronous design picks one of them, nearly
always the rising edge, and updates every flip-flop at that edge and only that one. Using a single
edge everywhere is what makes the timing of a large circuit analysable: every flip-flop changes at
the same well-defined instant, so you can reason one clock cycle at a time instead of tracking every
signal continuously.

Simulated and real clocks run on wildly different time scales, and the logic is identical in both
cases:
\begin{itemize}
  \item In CircuitVerse, set a period you can follow, e.g. \code{1000 ms}.
  \item On the DE0-CV the system clock is a fixed \code{50 MHz}, a period of \code{20 ns}. Nobody
    perceives individual edges at that rate, so you reason in terms of ``at the next rising edge'',
    never in real time.
\end{itemize}

% ------------------------------------------------------------------------------------------
\section{Registers: flip-flops in parallel}
\label{c3:sec:register}

A \textbf{register} stores more than one bit: N D flip-flops side by side, one per bit, all sharing
the same \code{clock}, \code{reset_n} and \code{enable}.

\begin{textdrawing}
D(0) ──►[D-FF]──► Q(0)
D(1) ──►[D-FF]──► Q(1)
D(2) ──►[D-FF]──► Q(2)
   ⋮        ⋮
D(N-1)──►[D-FF]──► Q(N-1)

  (clock, reset_n, enable are shared by every flip-flop above)
\end{textdrawing}

That is the whole idea: an 8-bit register is eight D flip-flops updating at the same clock edge.

In VHDL you rarely instantiate N flip-flops explicitly. Declare a \code{std_logic_vector} and
assign the whole of it inside a single synchronous process (\secref{c3:sec:template}): every bit
synthesises to a flip-flop of its own, and all of them share the process's clock and reset.

\begin{vhdlcode}
signal q: std_logic_vector(3 downto 0);
...
process (clock, reset_n) is
begin
    if (reset_n = '0') then
        q <= "0000";
    elsif (rising_edge(clock)) then
        if (enable = '1') then
            q <= d; -- d is a std_logic_vector(3 downto 0) input
        end if;
    end if;
end process;
\end{vhdlcode}

Four bits, one process, four flip-flops in hardware: in principle no different from four
\code{d_flip_flop} instances wired in parallel.

% ------------------------------------------------------------------------------------------
\section{The synchronous process template in VHDL}
\label{c3:sec:template}

Every synchronous element in this chapter, the flip-flop, the register and the edge detector, uses
the same shape. Without a reset, a single D flip-flop is:

\begin{vhdlcode}
process (clock) is
begin
    if (rising_edge(clock)) then
        q <= d;
    end if;
end process;
\end{vhdlcode}

\code{q} takes \code{d}'s value at every rising edge and holds it the rest of the time. That single
line inside the \code{if} statement is enough for synthesis to infer a real D flip-flop.

\subsection*{Rules}
\begin{itemize}
  \item \textbf{\code{clock} is the only signal in the sensitivity list}, unless there is an
    asynchronous reset. The process should run only when the clock changes; it has no reason to
    react to anything else.
  \item A process reacts to \emph{every} change of a signal in the sensitivity list, not just
    rising ones, so \code{rising_edge()} specifically filters out the rising edge. All the logic
    that updates synchronously sits inside that \code{if}.
  \item Every signal assigned inside such a process synthesises to a flip-flop, one per bit.
\end{itemize}

\subsection*{With an asynchronous reset}
An asynchronous reset must be in the sensitivity list too, since it can activate independently of
the clock, and it takes precedence when it is active.

\begin{vhdlcode}
process (clock, reset_n) is
begin
    if (reset_n = '0') then
        q <= '0';
    elsif (rising_edge(clock)) then
        q <= d;
    end if;
end process;
\end{vhdlcode}

Read it in priority order: check the reset first, and look for a rising edge only if it is
inactive. If neither holds (a falling edge, or \code{reset_n} going from \code{0} back to \code{1},
which also reruns the process) nothing is assigned and \code{q} keeps its previous value.

\subsection*{When the assignment actually takes effect}
One rule underlies all of the above, and it is the one place where \code{<=} behaves unlike
assignment in every language you have written software in.

A signal assignment does not take effect when its line runs. It \textbf{schedules} a value, which
is applied only once the process has finished this round and suspends. Until then, every read of
that signal gives the value it held \emph{before the round began}, however many assignments have
already run.

That sounds like a technicality. It is the reason the template above is a flip-flop and not a wire:

\begin{vhdlcode}
process (clock) is
begin
    if (rising_edge(clock)) then
        b <= a;
        c <= b;
    end if;
end process;
\end{vhdlcode}

Read as software, both \code{b} and \code{c} end up holding \code{a}. In hardware they do not. At
each edge \code{b} is scheduled to take \code{a}'s value, and \code{c} is scheduled to take the
value \code{b} held \emph{going into that edge}, not the one just scheduled onto it. The result is
two flip-flops in a chain: \code{a} reaches \code{b} after one edge and \code{c} after two. Swap
the lines and nothing changes, because neither assignment can observe the other's result.

Two consequences worth taking away:
\begin{itemize}
  \item \textbf{Order does not matter} between assignments to \emph{different} signals in the same
    clocked process. All of them read the old values and all of them take effect at once. That is
    the opposite of the sequential reading the word ``process'' invites, and it is why a chain of
    flip-flops can be written in any order.
  \item \textbf{Assign the same signal twice in one round and only the last assignment survives.}
    The earlier ones are discarded without ever reaching the signal.
\end{itemize}

\Secref{c3:sec:edge}'s edge detector is exactly the two-element chain above, and it works for
exactly this reason. \Chapref{c5:ch} returns to the rule from the other side, with a
\code{variable}, which updates immediately and therefore behaves the way the software reading
expects.

% ------------------------------------------------------------------------------------------
\section{Edge detection}
\label{c3:sec:edge}

The flip-flop's other everyday job is detecting the \emph{moment} a signal changes, rather than
reading its level: store the signal's previous value in a flip-flop and then compare it with the
current value using a gate.

\begin{itemize}
  \item \textbf{Rising edge detection} (\code{0} to \code{1}):
    $\mathit{edge} = \mathit{current} \cdot \mathit{previous}'$
  \item \textbf{Falling edge detection} (\code{1} to \code{0}):
    $\mathit{edge} = \mathit{current}' \cdot \mathit{previous}$
\end{itemize}

\begin{textdrawing}
signal ──┬───────────────────────────► AND ──► edge
         │                              ▲
         └──►[D-FF]── previous ──(NOT)──┘
                ▲
           clock, reset_n
\end{textdrawing}

\code{edge} is high for exactly one clock cycle, the one immediately after the transition, because
\code{previous} only catches up at the \emph{next} edge. That is the mechanism behind ``do
something once per button press'' instead of ``do something every clock cycle the button is held
down'', which is exactly what \secref{c3:sec:ledtoggle} solves.

% ------------------------------------------------------------------------------------------
\section{Worked example: a D flip-flop in VHDL}
\label{c3:sec:dffvhdl}

\file{d_flip_flop.vhd} implements \secref{c3:sec:dff} and \ref{c3:sec:template} together: an enable
gating whether a rising edge captures \code{d}, and an asynchronous active-low reset.

\vhdlfile{lectures/L03/d_flip_flop/d_flip_flop.vhd}

Its single process is the reset-plus-enable variant of \secref{c3:sec:template}'s template: the
reset takes precedence, then a rising edge captures \code{d} into an internal signal \code{q_s} if
\code{enable = '1'}. \code{q_s} exists because an \emph{output port} cannot be read back inside the
same architecture; \code{q} and \code{q_n} are both driven combinationally from it, outside the
process.

To try it on hardware, assign \code{clock} to the \code{50 MHz} system clock, \code{reset_n} to a
push button, \code{d} and \code{enable} to two slide switches, and \code{q}/\code{q_n} to two LEDs.
Then turn \code{enable} on and change \code{d}, and the LED follows; turn \code{enable} off and
change \code{d}, and it does not.

There is no clock edge you can perceive at \code{50 MHz}, so what you are really confirming is that
the logic still holds when you can no longer see the individual edges. The flip-flop appears to
respond to your switch; in fact it responded to one of fifty million clock edges that second, and
the switch only decided which value that edge captured.

% ------------------------------------------------------------------------------------------
\section{Worked example: edge-detected LED toggle}
\label{c3:sec:ledtoggle}

\file{led_toggle.vhd} combines \secref{c3:sec:register}, \ref{c3:sec:template} and
\ref{c3:sec:edge} into a small system: two independent push buttons, each toggling its own LED
exactly once per press.

\vhdlfile{lectures/L03/led_toggle/led_toggle.vhd}

Three 2-bit signals, two held in clocked processes and one computed between them:
\begin{itemize}
  \item \code{button_prev_n} is a 2-bit register (\secref{c3:sec:register}) holding last cycle's
    \code{button_n}, reset to \code{"11"} because the buttons are active low.
  \item \code{button_edge <= (not button_n) and button_prev_n} is \secref{c3:sec:edge}'s
    falling-edge detector applied to both bits at once, high for one cycle exactly when a button
    goes from released to pressed.
  \item \code{led_s} is a second 2-bit register which, instead of capturing a new value every
    cycle, \textbf{toggles} bit \code{i} only when \code{button_edge(i) = '1'}. Without that
    gating it would toggle every clock cycle instead of once per press, which is the whole point of
    computing \code{button_edge}.
\end{itemize}

\paragraph{On the board}
Create the Quartus Prime Lite project against the DE0-CV's \code{5CEBA4F23C7N} exactly as in
\chapref{c1:ch}, then assign \code{clock} to the \code{50 MHz} oscillator, \code{reset_n} and
\code{button_n(1 downto 0)} to three push buttons, and \code{led(1 downto 0)} to two LEDs. Every
press should toggle exactly one LED, once, however long you hold it down.

\begin{warning}
\textbf{One caveat, deliberately left open.} \code{button_n} is a raw signal straight from a
physical push-button pin. In simulation it is a clean, instantaneous transition; on real hardware it
is neither instantaneous (contacts bounce) nor synchronised to \code{clock}. This edge detector is
logically correct but \emph{unsafe} to wire directly to a real button on an FPGA without more work
first, and that work is exactly what \chapref{c4:ch} takes up.
\end{warning}
